\ifdefined\gkver\else\def\gkver{student}\fi
\documentclass[\gkver]{gaokaozhenti}
\begin{document}

\chapter{2006年江苏理综}
%% paperType: 理综物理部分

\begin{enumerate}
\item
%% number: 17
%% paperName: 2006年江苏理综第 17 题 3 分
%% typeId: 1
%% type: 单选题
%% chapter: 直线运动
%% point: 平均速度和瞬时速度
%% method: 无
%% score: 3
%% degree: 950
%% duplicateId: 0
%% body:
下列情况中的速度，属于平均速度的是\xzanswer{B}

\fourchoices[answer=B]
{百米赛跑的运动员冲过终点线时的速度为 $9.5\Ums$}
{由于堵车，汽车在通过隧道过程中的速度仅为 $1.2\Ums$}
{返回地球的太空舱落到太平洋水面时的速度为 $8\Ums$}
{子弹射到墙上时的速度为 $800\Ums$}

%% answer: B
%% memo:
\memoanswer{
本题原卷及材料中未提供详解，待补充。
}

\item
%% number: 18
%% paperName: 2006年江苏理综第 18 题 3 分
%% typeId: 1
%% type: 单选题
%% chapter: 力物体的平衡
%% point: 重力
%% method: 无
%% score: 3
%% degree: 850
%% duplicateId: 0
%% body:
关于重心，下列说法中正确的是\xzanswer{D}

\fourchoices[answer=D]
{重心就是物体内最重的一点}
{物体发生形变时，其重心位置一定不变}
{物体升高时，其重心在空中的位置一定不变}
{采用背越式跳高的运动员在越过横杆时，其重心位置可能在横杆之下}

%% answer: D
%% memo:
\memoanswer{
本题原卷及材料中未提供详解，待补充。
}

\item
%% number: 19
%% paperName: 2006年江苏理综第 19 题 3 分
%% typeId: 1
%% type: 单选题
%% chapter: 气体性质
%% point: 物体的内能
%% method: 无
%% score: 3
%% degree: 850
%% duplicateId: 0
%% body:
关于物体的内能，下列说法中正确的是\xzanswer{B}

\fourchoices[answer=B]
{物体的温度越高，其分子热运动的平均动能越小}
{物体的温度越高，其分子热运动的平均动能越大}
{只有做功才能改变物体的内能}
{只有热传递才能改变物体的内能}

%% answer: B
%% memo:
\memoanswer{
本题原卷及材料中未提供详解，待补充。
}

\item
%% number: 20
%% paperName: 2006年江苏理综第 20 题 3 分
%% typeId: 1
%% type: 单选题
%% chapter: 电场
%% point: 库仑定律
%% method: 无
%% score: 3
%% degree: 850
%% duplicateId: 0
%% body:
如图所示，用两根同样的绝缘细线把甲、乙两个质量相等的带电小球悬挂在同一点上，甲、乙两球均处于静止状态。已知两球带同种电荷，且甲球的电荷量大于乙球的电荷量，$F_{1}$、$F_{2}$ 分别表示甲、乙两球所受的库仑力，则下列说法中正确的是\xzanswer{C}

\onepicture[w=4.0cm]{figs/20.png}

\fourchoices[answer=C]
{$F_{1}$ 一定大于 $F_{2}$}
{$F_{1}$ 一定小于 $F_{2}$}
{$F_{1}$ 与 $F_{2}$ 大小一定相等}
{无法比较 $F_{1}$ 与 $F_{2}$ 的大小}

%% answer: C
%% memo:
\memoanswer{
本题原卷及材料中未提供详解，待补充。
}

\item
%% number: 21
%% paperName: 2006年江苏理综第 21 题 3 分
%% typeId: 1
%% type: 单选题
%% chapter: 磁场
%% point: 磁感线
%% method: 无
%% score: 3
%% degree: 950
%% duplicateId: 0
%% body:
关于磁感线，下列说法中正确的是\xzanswer{B}

\fourchoices[answer=B]
{磁感线是实际存在于磁场中的线}
{磁感线上任意一点的切线方向，都跟该点的磁场方向一致}
{磁感线是一条条不闭合的曲线}
{磁感线有可能出现相交的情况}

%% answer: B
%% memo:
\memoanswer{
本题原卷及材料中未提供详解，待补充。
}

\item
%% number: 22
%% paperName: 2006年江苏理综第 22 题 3 分
%% typeId: 1
%% type: 单选题
%% chapter: 原子和原子核
%% point: 原子核的组成
%% method: 无
%% score: 3
%% degree: 950
%% duplicateId: 0
%% body:
下列说法中正确的是\xzanswer{A}

\fourchoices[answer=A]
{原子核由质子和中子组成}
{原子核由质子和电子组成}
{质子和中子都带正电}
{原子核的质量数一定等于电荷数}

%% answer: A
%% memo:
\memoanswer{
本题原卷及材料中未提供详解，待补充。
}

\item
%% number: 23
%% paperName: 2006年江苏理综第 23 题 3 分
%% typeId: 1
%% type: 单选题
%% chapter: 曲线运动
%% point: 人造地球卫星
%% method: 无
%% score: 3
%% degree: 750
%% duplicateId: 0
%% body:
举世瞩目的“神舟”六号航天飞船的成功发射和顺利返回，显示了我国航天事业取得的巨大成就。已知地球的质量为 $M$，引力常量为 $G$，设飞船绕地球做匀速圆周运动的轨道半径为 $r$，则飞船在圆轨道上运行的速率为\xzanswer{A}

\fourchoices[answer=A]
{$\sqrt{\frac{GM}{r}}$}
{$\sqrt{\frac{r}{GM}}$}
{$\sqrt{\frac{G}{Mr}}$}
{$\sqrt{\frac{M}{Gr}}$}

%% answer: A
%% memo:
\memoanswer{
本题原卷及材料中未提供详解，待补充。
}

\item
%% number: 24
%% paperName: 2006年江苏理综第 24 题 3 分
%% typeId: 1
%% type: 单选题
%% chapter: 光学
%% point: 全反射
%% method: 无
%% score: 3
%% degree: 850
%% duplicateId: 0
%% body:
在医学上，光导纤维可以制成内窥镜，用来检查人体胃、肠、气管等器官的内部。内窥镜有两组光导纤维，一组用来把光输送到人体内部，另一组用来进行观察。光在光导纤维中的传输利用了\xzanswer{D}

\fourchoices[answer=D]
{光的折射}
{光的衍射}
{光的干涉}
{光的全反射}

%% answer: D
%% memo:
\memoanswer{
本题原卷及材料中未提供详解，待补充。
}

\item
%% number: 25
%% paperName: 2006年江苏理综第 25 题 3 分
%% typeId: 1
%% type: 单选题
%% chapter: 电路
%% point: 多用表的使用
%% method: 无
%% score: 3
%% degree: 850
%% duplicateId: 0
%% body:
关于多用电表的使用，下列说法中正确的是\xzanswer{A}

\fourchoices[answer=A]
{用电流挡测电流或用电压挡测电压前，必须检查机械零点}
{用电阻挡测电阻前，不需要检查机械零点}
{用电阻挡测电阻时，若从一个倍率变换到另一个倍率，不需要重新检查欧姆零点}
{用电阻挡测电阻时，被测电阻的阻值越大，指针向右转过的角度就越大}

%% answer: A
%% memo:
\memoanswer{
本题原卷及材料中未提供详解，待补充。
}

\item
%% number: 32
%% paperName: 2006年江苏理综第 32 题 4 分
%% typeId: 1
%% type: 单选题
%% chapter: 机械振动
%% point: 振动图象
%% method: 无
%% score: 4
%% degree: 950
%% duplicateId: 0
%% body:
如图所示为一个弹簧振子做简谐运动的振动图象，由图可知，该弹簧振子的振幅是 \tkanswer[1.5cm]{10}$\Ucm$，振动频率是 \tkanswer[1.5cm]{0.5}$\UHz$。

\onepicture[w=7.0cm]{figs/32.png}

%% answer: 
%% memo:
\memoanswer{
本题原卷及材料中未提供详解，待补充。
}

\item
%% number: 33
%% paperName: 2006年江苏理综第 33 题 4 分
%% typeId: 3
%% type: 填空题
%% chapter: 电磁感应
%% point: 法拉第电磁感应定律
%% method: 无
%% score: 4
%% degree: 750
%% duplicateId: 0
%% body:
一个 500 匝的线圈，其电阻为 $5\UO$，将它与电阻为 $495\UO$ 的电热器连成闭合电路。若在 $0.3\Us$ 内，穿过线圈的磁通量从 $0.03\UWb$ 均匀增加到 $0.09\UWb$，则线圈中产生的感应电动势为 \tkanswer[1.5cm]{100}$\UV$，通过电热器的电流为 \tkanswer[1.5cm]{0.2}$\UA$。

%% answer: 
\jdanswer{
100，0.2
}

%% memo:
\memoanswer{
本题原卷及材料中未提供详解，待补充。
}

\item
%% number: 34
%% paperName: 2006年江苏理综第 34 题 7 分
%% typeId: 6
%% type: 计算题
%% chapter: 运动和力
%% point: 牛顿第二定律
%% method: 无
%% score: 7
%% degree: 750
%% duplicateId: 0
%% body:
列车在机车的牵引下沿平直铁轨匀加速行驶，在 $100\Us$ 内速度由 $5.0\Ums$ 增加到 $15.0\Ums$。

\begin{enumerate}
	\item
	求列车的加速度大小；

	\item
	若列车的质量是 $1.0\times10^{6}\Ukg$，机车对列车的牵引力是 $1.5\times10^{5}\UN$，求列车在运动中所受的阻力大小。
\end{enumerate}

%% answer: 
\jdanswer{
\begin{enumerate}
	\item
	$a=0.1\Umsq$；

	\item
	$f=5.0\times10^{4}\UN$。

\end{enumerate}
}

%% memo:
\memoanswer{
（1）根据 $a=\frac{v_{t}-v_{0}}{t}$，代入数据得 $a=0.1\Umsq$。

（2）设列车在运动中所受的阻力大小为 $f$，由牛顿第二定律 $F_{\text{合}}=F_{\text{牵}}-f=ma$，代入数据解得 $f=5.0\times10^{4}\UN$。

评分标准：本题共 7 分，其中第（1）问 3 分，第（2）问 4 分。
}

\item
%% number: 35
%% paperName: 2006年江苏理综第 35 题 9 分
%% typeId: 1
%% type: 单选题
%% chapter: 电路
%% point: 闭合电路欧姆定律
%% method: 无
%% score: 9
%% degree: 650
%% duplicateId: 0
%% body:
如图所示电路中，电阻 $R_{1}=R_{2}=R_{3}=10\UO$，电源内阻 $r=5\UO$，电压表可视为理想电表。当开关 $S_{1}$ 和 $S_{2}$ 均闭合时，电压表的示数为 $10\UV$。

\begin{enumerate}
	\item
	电阻 $R_{2}$ 中的电流为多大？

	\item
	路端电压为多大？

	\item
	电源的电动势为多大？

	\item
	当开关 $S_{1}$ 闭合而 $S_{2}$ 断开时，电压表的示数变为多大？
\end{enumerate}

\onepicture[w=7.0cm, align=right]{figs/35.png}

%% answer: IAE
%% memo:
\memoanswer{
（1）电阻 $R_{2}$ 中的电流 $I=\frac{U_{2}}{R_{2}}$，代入数据得 $I=1\UA$。

（2）外电阻 $R=R_{2}+\frac{R_{1}R_{3}}{R_{1}+R_{3}}=15\UO$，路端电压 $U=IR=15\UV$。

（3）根据闭合电路欧姆定律 $I=\frac{E}{R+r}$，代入数据解得 $E=20\UV$。

（4）$S_{1}$ 闭合而 $S_{2}$ 断开，电路中的总电流 $I'=\frac{E}{R_{1}+R_{2}+r}$，电压表示数 $U'=I'(R_{1}+R_{2})$，代入数据解得 $U'=16\UV$。

评分标准：本题共 9 分，其中第（1）～（3）问各 2 分，第（4）问 3 分。
}

\item
%% number: 36
%% paperName: 2006年江苏理综第 36 题 4 分
%% typeId: 3
%% type: 填空题
%% chapter: 功和能
%% point: 机械能守恒定律
%% method: 无
%% score: 4
%% degree: 650
%% duplicateId: 0
%% body:
水力发电利用了水能这一可再生能源。在西部大开发的壮丽画卷中，三峡工程是浓墨重彩的一笔。设三峡水库水面到发电站水轮机的落差为 $h$，重力加速度为 $g$，不计水的初速和所受的阻力，则水到达水轮机时的速率等于 \tkanswer[2.0cm]{$\sqrt{2gh}$}。若每秒有质量为 $m$ 的水通过水轮机，水力发电中将水的动能转化为电能的效率为 $\eta$，则发电的功率等于 \tkanswer[2.0cm]{$\eta mgh$}。

%% answer: 
\jdanswer{
$\sqrt{2gh}$，$\eta mgh$
}

%% memo:
\memoanswer{
本题原卷及材料中未提供详解，待补充。
}

\end{enumerate}

\end{document}
